\documentclass[11pt,a4paper]{article}
\usepackage[margin=0.85in,top=1in,bottom=1in]{geometry}
\usepackage{amsmath,amssymb,amsfonts}
\usepackage{graphicx}
\usepackage{float}
\usepackage{enumitem}
\usepackage{microtype}
\usepackage{caption}
\usepackage[hidelinks]{hyperref}
\usepackage[most,breakable]{tcolorbox}
\tcbuselibrary{breakable,skins}
\usepackage[utf8]{inputenc}
\usepackage[T1]{fontenc}
\usepackage{lmodern}
\captionsetup{font=small,labelfont=bf,skip=4pt}
\setlist{leftmargin=1.5em,topsep=2pt}

%% Breakable card — prevents content cutoff across pages
\newtcolorbox{solcard}[3]{
  breakable, enhanced,
  colback=white, colframe=gray!50, arc=2pt,
  title={\textbf{#1}\hfill\textsf{\small #3\,pts}},
  fonttitle=\normalsize\sffamily\bfseries,
  before upper={\small\textit{#2}\par\smallskip\hrule height 0.4pt\smallskip},
  top=4pt, bottom=6pt, left=7pt, right=7pt,
  parbox=false,
}

\title{\textbf{Anisotropy of thermal conductivity --- Solution}}
\author{\normalsize X23 (2023) — English translation}
\date{}
\begin{document}\maketitle\vspace{-0.5em}
\begin{solcard}{A1}{Using the least squares method, obtain formulas for finding the coefficients $A$ and $B$ for each temperature value $T_0$ and for each material. Express your answer in terms of sample average values of products of the form $\Delta T^{k_1}d^{k_2}$, where $k_{1,2}$ are some numbers.}{0.60}

According to the method described in the condition, we compose a system of equations:\[\begin{cases} \cfrac{\partial S}{\partial A}=0\\ \cfrac{\partial S}{\partial B}=0\end{cases}, \quad  \text{where} \quad\quad S=\sum_i\left(\Delta T_{i\ \mathrm{exp}}-\Delta T_{i\ \mathrm{th}}(d_i, A, B)\right)^2=\sum_i\left(\Delta T_i-Ad_i+BT_0^{1/2}d_i^{3/2}\right)^2\]Thus, the system will take the form:\[\begin{cases} \cfrac\partial{\partial A}\sum_i\left(\Delta T_i-Ad_i+BT_0^{1/2}d_i^{3/2}\right)^2=0\\\cfrac\partial{\partial B}\sum_i\left(\Delta T_i-Ad_i+BT_0^{1/2}d_i^{3/2}\right)^2=0\end{cases}\implies\begin{cases}0=\sum_i\left(\Delta T_id_i-Ad_i^2+BT_0^{1/2}d_i^{5/2}\right)\\0=\sum_i\left(\Delta T_id_i^{3/2}-Ad_i^{5/2}+BT_0^{1/2}d_i^3\right)\end{cases}\]This is a system of linear equations for the coefficients $A$ and $B$. The coefficients in it are necessary products of the form $\Delta T^{k_1}d^{k_2}$. Using the average values of these products over the experimental sample, this system is written somewhat more elegantly in the form:\[\begin{cases}0=\overline{\Delta Td}-A\overline{d^2}+BT_0^{1/2}\overline{d^{5/2}}\\0=\overline{\Delta Td^{3/2}}-A\overline{d^{5/2}}+BT_0^{1/2}\overline{d^3}\end{cases}\]Solving this system, we finally obtain:\[\begin{array}{l}A=\cfrac{\overline{\Delta Td^{3/2}}\cdot\overline{d^{5/2}}-\overline{\Delta Td}\cdot\overline{d^3}}{\overline{d^{5/2}}^2-\overline{d^2}\cdot\overline{d^3}}\\B=\cfrac1{\sqrt{T_0}}\cfrac{\overline{\Delta Td^{3/2}}\cdot\overline{d^2}-\overline{\Delta Td}\cdot\overline{d^{5/2}}}{\overline{d^{5/2}}^2-\overline{d^2}\cdot\overline{d^3}}\end{array}\]

\par\smallskip\noindent\textbf{Answer:} \[\begin{array}{l}A=\cfrac{\overline{\Delta Td^{3/2}}\cdot\overline{d^{5/2}}-\overline{\Delta Td}\cdot\overline{d^3}}{\overline{d^{5/2}}^2-\overline{d^2}\cdot\overline{d^3}}\\B=\cfrac1{\sqrt{T_0}}\cfrac{\overline{\Delta Td^{3/2}}\cdot\overline{d^2}-\overline{\Delta Td}\cdot\overline{d^{5/2}}}{\overline{d^{5/2}}^2-\overline{d^2}\cdot\overline{d^3}}\end{array}\]
\end{solcard}

\begin{solcard}{A2}{For molybdenum sulfide $\rm(MoS_2)$, find $A$ and $B$ for each of the temperatures $T_0$.}{1.20}

$T_0=273~\mathrm{K}$. $\overline{d},~\mu\text{m}$ 20 40 60 80 100 120 140 160 180 200 $\Delta T,~\mathrm{K}$ 0.138 0.266 0.394 0.532 0.662 0.777 0.831 0.954 1.137 1.189 0.688 $\Delta Td,~\mathrm{K}\cdot\mu\text{m}$ 2.76 10.64 23.64 42.56 66.20 93.24 116.34 152.64 204.66 237.80 95.05 $\Delta Td^{3/2},~\mathrm{K}\cdot\mu\text{m}^{3/2}$ 12.34 6729 18311 380.67 662.00 1021.39 1376.55 1930.76 2745.80 3363.00 1174.29 $d^2,~\mathrm{\mu m}^2$ 400 1600 3600 6400 10000 14400 19600 25600 32400 40000 15400 $d^{5/2},~\mathrm{\mu m}^{5/2}$ 1788.85 10119.29 27885.48 57243.34 100000.00 157744.10 231910.33 323817.23 434691.61 565685.42 191088.57 $d^3,~\mathrm{\mu m}^3$ 8000 64000 216000 512000 1000000 1728000 2744000 4096000 5832000 8000000 2420000 $T_0=323~\mathrm{K}\overline{d},~\mathrm{\mu m}$ 20 40 60 80 100 120 140 160 180 200 $\Delta T,~\mathrm{K}$ 0.137 0.271 0.390 0.496 0.577 0.719 0.788 0.877 1.046 1.098 0.640 $\Delta Td,~\mathrm{K}\cdot\mu\text{m}$ 2.74 10.84 23.40 39.68 57.70 86.28 110.32 140.32 188.28 219.60 87.92 $\Delta Td^{3/2},~\mathrm{K}\cdot\mu\text{m}^{3/2}$ 12.25 68.56 181.26 354.91 577.00 945.15 1305.32 1774.92 2526.04 3105.61 1085.10 $d^2,~\mathrm{\mu m}^2$ 400 1600 3600 6400 10000 14400 19600 25600 32400 40000 15400 $d^{5/2},~\mathrm{\mu m}^{5/2}$ 1788.85 10119.29 27885.48 57243.34 100000.00 157744.10 231910.33 323817.23 434691.61 565685.42 191088.57 $d^3,~\mathrm{\mu m}^3$ 8000 64000 216000 512000 1000000 1728000 2744000 4096000 5832000 8000000 2420000 $T_0=373~\mathrm{K}\overline{d},~\mathrm{\mu m}$ 20 40 60 80 100 120 140 160 180 200 $\Delta T,~\mathrm{K}$ 0.139 0.267 0.373 0.483 0.563 0.676 0.760 0.856 0.901 1.023 0.604 $\Delta Td,~\mathrm{K}\cdot\mu\text{m}$ 2.78 10.68 22.38 38.64 56.30 81.12 106.40 136.96 162.18 204.60 82.20 $\Delta Td^{3/2},~\mathrm{K}\cdot\mu\text{m}^{3/2}$ 12.43 67.55 173.35 345.61 563.00 888.63 1258.94 1732.42 2175.87 2893.48 1011.13 $d^2,~\mathrm{\mu m}^2$ 400 1600 3600 6400 10000 14400 19600 25600 32400 40000 15400 $d^{5/2},~\mathrm{\mu m}^{5/2}$ 1788.85 10119.29 27885.48 57243.34 100000.00 157744.10 231910.33 323817.23 434691.61 565685.42 191088.57 $d^3,~\mathrm{\mu m}^3$ 8000 64000 216000 512000 1000000 1728000 2744000 4096000 5832000 8000000 2420000

\par\smallskip\noindent\textbf{Answer:} $T_0,~\mathrm{K}$ 273 323 373 $A,~\cfrac{\mathrm{K}}{\mathrm{mm}}$ 7.46 7.18 7.59 $B,~\cfrac{\mathrm{K}^{1/2}}{\mathrm{mm}^{3/2}}$ 0.199 0.208 0.298
\end{solcard}

\begin{solcard}{A3}{Find out whether it is possible to conclude from the results of the experiment that thermal conductivity $\varkappa_\perp$ depends on temperature, underline the correct option on the answer sheet (both sides of the inequality must be explicitly calculated in the solution). Find the arithmetic mean $\overline A$ and calculate the thermal conductivity $\varkappa_\perp$ from it.}{0.50}

We substitute the numbers (taking into account that $T_{0\max}=373~\mathrm{K}$, $d_{\max}=0.2~\mathrm{mm}$): \[\Delta A < \overline BT_{0\max}^{1/2}d_{\max}^{1/2}\implies0.41~\cfrac{\mathrm{K}}{\mathrm{mm}} < 2.03~\cfrac{\mathrm{K}}{\mathrm{mm}}~-~\text{depends}~/~\underline{no~\text{depends}}.\]

\par\smallskip\noindent\textbf{Answer:} \[\overline{A}=7.41~\cfrac{\mathrm{K}}{\mathrm{mm}}\\\ \\\varkappa_\perp=53.3~\cfrac{\mathrm{mW}}{\text{m}\cdot\text{K}}\ne\varkappa_\perp(T_0)\]
\end{solcard}

\begin{solcard}{A4}{Follow steps $\bf A2$ and $\bf A3$ for tungsten sulfide $\rm(WS_2)$.}{1.70}

$T_0=273~\mathrm{K}$. $\overline{d},~\mu\text{m}$ 20 40 60 80 100 120 140 160 180 200 $\Delta T,~\mathrm{K}$ 0.191 0.381 0.571 0.764 0.946 1.100 1.304 1.443 1.577 1.809 1.009 $\Delta Td,~\mathrm{K}\cdot\mu\text{m}$ 3.82 15.24 34.26 61.12 94.60 132.00 182.56 230.88 283.86 361.80 140.01 $\Delta Td^{3/2},~\mathrm{K}\cdot\mu\text{m}^{3/2}$ 17.08 96.39 265.38 546.67 946.00 1445.99 2160.08 2920.43 3808.38 5116.62 1732.30 $d^2,~\mathrm{\mu m}^2$ 400 1600 3600 6400 10000 14400 19600 25600 32400 40000 15400 $d^{5/2},~\mathrm{\mu m}^{5/2}$ 1788.85 10119.29 27885.48 57243.34 100000.00 157744.10 231910.33 323817.23 434691.61 565685.42 191088.57 $d^3,~\mathrm{\mu m}^3$ 8000 64000 216000 512000 1000000 1728000 2744000 4096000 5832000 8000000 2420000 $T_0=323~\mathrm{K}\overline{d},~\mathrm{\mu m}$ 20 40 60 80 100 120 140 160 180 200 $\Delta T,~\mathrm{K}$ 0.195 0.384 0.566 0.727 0.907 1.083 1.202 1.354 1.518 1.701 0.964 $\Delta Td,~\mathrm{K}\cdot\mu\text{m}$ 3.9 15.36 33.96 58.16 90.7 129.96 168.28 216.64 273.24 340.2 133.04 $\Delta Td^{3/2},~\mathrm{K}\cdot\mu\text{m}^{3/2}$ 17.44 97.15 263.05 520.20 907.00 1423.64 1991.12 2740.30 3665.90 4811.15 1643.70 $d^2,~\mathrm{\mu m}^2$ 400 1600 3600 6400 10000 14400 19600 25600 32400 40000 15400 $d^{5/2},~\mathrm{\mu m}^{5/2}$ 1788.85 10119.29 27885.48 57243.34 100000.00 157744.10 231910.33 323817.23 434691.61 565685.42 191088.57 $d^3,~\mathrm{\mu m}^3$ 8000 64000 216000 512000 1000000 1728000 2744000 4096000 5832000 8000000 2420000 $T_0=373~\mathrm{K}\overline{d},~\mathrm{\mu m}$ 20 40 60 80 100 120 140 160 180 200 $\Delta T,~\mathrm{K}$ 0.201 0.367 0.549 0.733 0.877 1.037 1.172 1.330 1.439 1.629 0.933 $\Delta Td,~\mathrm{K}\cdot\mu\text{m}$ 4.02 14.68 32.94 58.64 87.70 124.44 164.08 212.80 259.02 325.80 128.41 $\Delta Td^{3/2},~\mathrm{K}\cdot\mu\text{m}^{3/2}$ 17.98 92.84 255.15 524.49 877.00 1363.17 1941.42 2691.73 3475.12 4607.51 1584.64 $d^2,~\mathrm{\mu m}^2$ 400 1600 3600 6400 10000 14400 19600 25600 32400 40000 15400 $d^{5/2},~\mathrm{\mu m}^{5/2}$ 1788.85 10119.29 27885.48 57243.34 100000.00 157744.10 231910.33 323817.23 434691.61 565685.42 191088.57 $d^3,~\mathrm{\mu m}^3$ 8000 64000 216000 512000 1000000 1728000 2744000 4096000 5832000 8000000 2420000 Substitute the numbers (taking into account that $T_{0\max}=373~\mathrm{K}$, $d_{\max}=0.2~\mathrm{mm}$): \[\Delta A < \overline BT_{0\max}^{1/2}d_{\max}^{1/2}\implies0.23~\cfrac{\mathrm{K}}{\mathrm{mm}} < 2.17~\cfrac{\mathrm{K}}{\mathrm{mm}}~-~\text{depends}~/~\underline{no~\text{depends}}\]


We substitute the numbers (taking into account that $T_{0\max}=373~\mathrm{K}$, $d_{\max}=0.2~\mathrm{mm}$): \[\Delta A < \overline BT_{0\max}^{1/2}d_{\max}^{1/2}\implies0.23~\cfrac{\mathrm{K}}{\mathrm{mm}} < 2.17~\cfrac{\mathrm{K}}{\mathrm{mm}}~-~\text{depends}~/~\underline{no~\text{depends}}\]


\[\overline{A}=10.47~\cfrac{\mathrm{K}}{\mathrm{mm}}\\\varkappa_\perp=37.7~\cfrac{\mathrm{mW}}{\text{m}\cdot\text{K}}\ne\varkappa_\perp(T_0)\]

\par\smallskip\noindent\textbf{Answer:} $T_0,~\mathrm{K}$ 273 323 373 $A,~\cfrac{\mathrm{K}}{\mathrm{mm}}$ 10.37 10.44 10.60 $B,~\cfrac{\mathrm{K}^{1/2}}{\mathrm{mm}^{3/2}}$ 0.197 0.256 0.300 \[\overline{A}=10.47~\cfrac{\mathrm{K}}{\mathrm{mm}}\\\varkappa_\perp=37.7~\cfrac{\mathrm{mW}}{\text{m}\cdot\text{K}}\ne\varkappa_\perp(T_0)\]
\end{solcard}

\begin{solcard}{B1}{Using the least squares method, obtain a formula for finding the coefficient $A$ for each temperature value $T_0$ and for each material. Express your answer in terms of sample average values of products of the form $\Delta T^{k_1}d^{k_2}$, where $k_{1,2}$ are some numbers.}{0.60}

According to the method described in the condition, we compose a system of equations:\[\begin{cases} \cfrac{\partial S}{\partial A}=0\\ \cfrac{\partial S}{\partial B}=0\end{cases}, \quad  \text{where} \quad\quad S=\sum_i\left(\Delta T_{i\ \mathrm{exp}}-\Delta T_{i\ \mathrm{th}}(d_i, A, B)\right)^2=\sum_i\left(\Delta T_i-\cfrac A{d_i}+\cfrac B{d_i^2}\right)^2\]Thus, the system will take the form:\[\begin{cases} \cfrac\partial{\partial A}\sum_i\left(\Delta T_i-\cfrac A{d_i}+\cfrac B{d_i^2}\right)^2=0\\\cfrac\partial{\partial B}\sum_i\left(\Delta T_i-\cfrac A{d_i}+\cfrac B{d_i^2}\right)^2=0\end{cases}\implies\begin{cases}0=\sum_i\left(\Delta T_id_i^{-1}-Ad_i^{-2}+Bd_i^{-3}\right)\\0=\sum_i\left(\Delta T_id_i^{-2}-Ad_i^{-3}+Bd_i^{-4}\right)\end{cases}\]This is a system of linear equations for the coefficients $A$ and $B$. The coefficients in it are necessary products of the form $\Delta T^{k_1}d^{k_2}$. Using the average values of these products over the experimental sample, this system is written somewhat more elegantly in the form:\[\begin{cases}0=\overline{\Delta Td^{-1}}-A\overline{d^{-2}}+B\overline{d^{-3}}\\0=\overline{\Delta Td^{-2}}-A\overline{d^{-3}}+B\overline{d^{-4}}\end{cases}\]Solving this system for $A$, we finally obtain:\[A=\frac{\overline{\Delta Td^{-2}}\cdot\overline{d^{-3}}-\overline{\Delta Td^{-1}}\cdot\overline{d^{-4}}}{\overline{d^{-3}}^2-\overline{d^{-2}}\cdot\overline{d^{-4}}}.\]

\par\smallskip\noindent\textbf{Answer:} \[A=\frac{\overline{\Delta Td^{-2}}\cdot\overline{d^{-3}}-\overline{\Delta Td^{-1}}\cdot\overline{d^{-4}}}{\overline{d^{-3}}^2-\overline{d^{-2}}\cdot\overline{d^{-4}}}\]
\end{solcard}

\begin{solcard}{B2}{For molybdenum sulfide $\rm(MoS_2)$, find $A$ for each of the temperatures $T_0$. From here, for each $T_0$, find the values of $\varkappa_\parallel$. Find the thermal conductivity anisotropy coefficient $\rho=\cfrac{\varkappa_\parallel}{\varkappa_\perp}$ at temperature $T_0=273\ \mathrm{K}$.}{1.00}

$T_0=273~\mathrm{K}$. $\overline{d},~\mu\text{m}$ 120 140 160 180 200 $\Delta T,~\mathrm{K}$ 1.372 1.157 1.055 0.948 0.846 $\Delta Td^{-1},~10^{-3}~\mathrm{K}\cdot\mu\text{m}^{-1}$ 11.43 8.26 6.59 5.27 4.23 7.16 $\Delta Td^{-2},~10^{-6}~\mathrm{K}\cdot\mu\text{m}^{-2}$ 95.28 59.03 41.21 29.26 21.15 49.19 $d^{-2},~10^{-6}~\mathrm{\mu m}^{-2}$ 69.44 51.02 39.06 30.86 25.00 43.08 $d^{-3},~10^{-9}~\mathrm{\mu m}^{-3}$ 578.70 364.43 244.14 171.47 125.00 296.75 $d^{-4},~10^{-9}~\mathrm{\mu m}^{-4}$ 4.823 2.603 1.526 0.953 0.625 2.106 $T_0=323~\mathrm{K}\overline{d},~\mathrm{\mu m}$ 120 140 160 180 200 $\Delta T,~\mathrm{K}$ 1.407 1.247 1.144 1.019 0.948 $\Delta Td^{-1},~10^{-3}~\mathrm{K}\cdot\mu\text{m}^{-1}$ 11.73 8.91 7.15 5.66 4.74 7.64 $\Delta Td^{-2},~10^{-6}~\mathrm{K}\cdot\mu\text{m}^{-2}$ 97.71 63.62 44.69 31.45 23.70 52.23 $d^{-2},~10^{-6}~\mathrm{\mu m}^{-2}$ 69.44 51.02 39.06 30.86 25.00 43.08 $d^{-3},~10^{-9}~\mathrm{\mu m}^{-3}$ 578.70 364.43 244.14 171.47 125.00 296.75 $d^{-4},~10^{-9}~\mathrm{\mu m}^{-4}$ 4.823 2.603 1.526 0.953 0.625 2.106 $T_0=373~\mathrm{K}\overline{d},~\mathrm{\mu m}$ 120 140 160 180 200 $\Delta T,~\mathrm{K}$ 1.455 1.356 1.194 1.107 1.024 $\Delta Td^{-1},~10^{-3}~\mathrm{K}\cdot\mu\text{m}^{-1}$ 12.13 9.69 7.46 6.15 5.12 8.11 $\Delta Td^{-2},~10^{-6}~\mathrm{K}\cdot\mu\text{m}^{-2}$ 101.04 69.18 46.64 34.17 25.60 55.33 $d^{-2},~10^{-6}~\mathrm{\mu m}^{-2}$ 69.44 51.02 39.06 30.86 25.00 43.08 $d^{-3},~10^{-9}~\mathrm{\mu m}^{-3}$ 578.70 364.43 244.14 171.47 125.00 296.75 $d^{-4},~10^{-9}~\mathrm{\mu m}^{-4}$ 4.823 2.603 1.526 0.953 0.625 2.106


\[\rho=9.85\cdot10^2\]

\par\smallskip\noindent\textbf{Answer:} $T_0,~\mathrm{K}$ 273 323 373 $A,~\mathrm{\mu m}\cdot\text{K}$ 179.6 218.8 247.5 $\varkappa_\parallel,~\cfrac{\text{W}}{\text{m}\cdot\text{K}}$ 52.5 43.1 38.1 \[\rho=9.85\cdot10^2\]
\end{solcard}

\begin{solcard}{B3}{Find $n$.}{0.40}

You can calculate coefficients $\varkappa_{\parallel,0}$ and $n$ using least squares for a modeling function of the form: \[\varkappa_\parallel=\frac{\varkappa_{\parallel,0}}{T_0^n}\] using a standard calculator.

\par\smallskip\noindent\textbf{Answer:} \[n=1.03\]
\end{solcard}

\begin{solcard}{B4}{Follow steps $\bf B2$ and $\bf B3$ for tungsten sulfide $\rm(WS_2)$.}{1.40}

$T_0=273~\mathrm{K}$. $\overline{d},~\mu\text{m}$ 120 140 160 180 200 160 $\Delta T,~\mathrm{K}$ 1.596 1.417 1.284 1.163 1.084 1.309 $\Delta Td^{-1},~10^{-3}~\mathrm{K}\cdot\mu\text{m}^{-1}$ 13.30 10.12 8.03 6.46 5.42 8.67 $\Delta Td^{-2},~10^{-6}~\mathrm{K}\cdot\mu\text{m}^{-2}$ 110.83 72.30 50.16 35.90 27.10 59.26 $d^{-2},~10^{-6}~\mathrm{\mu m}^{-2}$ 69.44 51.02 39.06 30.86 25.00 43.08 $d^{-3},~10^{-9}~\mathrm{\mu m}^{-3}$ 578.70 364.43 244.14 171.47 125.00 296.75 $d^{-4},~10^{-9}~\mathrm{\mu m}^{-4}$ 4.823 2.603 1.526 0.953 0.625 2.106 $T_0=323~\mathrm{K}\overline{d},~\mathrm{\mu m}$ 120 140 160 180 200 160 $\Delta T,~\mathrm{K}$ 1.678 1.482 1.370 1.255 1.162 1.389 $\Delta Td^{-1},~10^{-3}~\mathrm{K}\cdot\mu\text{m}^{-1}$ 13.98 10.59 8.56 6.97 5.81 9.18 $\Delta Td^{-2},~10^{-6}~\mathrm{K}\cdot\mu\text{m}^{-2}$ 116.53 75.61 53.52 38.73 29.05 62.69 $d^{-2},~10^{-6}~\mathrm{\mu m}^{-2}$ 69.44 51.02 39.06 30.86 25.00 43.08 $d^{-3},~10^{-9}~\mathrm{\mu m}^{-3}$ 578.70 364.43 244.14 171.47 125.00 296.75 $d^{-4},~10^{-9}~\mathrm{\mu m}^{-4}$ 4.823 2.603 1.526 0.953 0.625 2.106 $T_0=373~\mathrm{K}\overline{d},~\mathrm{\mu m}$ 120 140 160 180 200 160 $\Delta T,~\mathrm{K}$ 1.763 1.607 1.509 1.355 1.307 1.508 $\Delta Td^{-1},~10^{-3}~\mathrm{K}\cdot\mu\text{m}^{-1}$ 14.69 11.48 9.43 7.53 6.54 9.93 $\Delta Td^{-2},~10^{-6}~\mathrm{K}\cdot\mu\text{m}^{-2}$ 122.43 81.99 58.95 41.82 32.68 67.57 $d^{-2},~10^{-6}~\mathrm{\mu m}^{-2}$ 69.44 51.02 39.06 30.86 25.00 43.08 $d^{-3},~10^{-9}~\mathrm{\mu m}^{-3}$ 578.70 364.43 244.14 171.47 125.00 296.75 $d^{-4},~10^{-9}~\mathrm{\mu m}^{-4}$ 4.823 2.603 1.526 0.953 0.625 2.106 You can calculate coefficients $\varkappa_{\parallel,0}$ and $n$ using least squares for a modeling function of the form: \[\varkappa_\parallel=\frac{\varkappa_{\parallel,0}}{T_0^n}\] using a standard calculator.


You can calculate coefficients $\varkappa_{\parallel,0}$ and $n$ using least squares for a modeling function of the form: \[\varkappa_\parallel=\frac{\varkappa_{\parallel,0}}{T_0^n}\] using a standard calculator.


\[\rho=1.000\cdot10^3\\n=1.02\]

\par\smallskip\noindent\textbf{Answer:} $T_0,~\mathrm{K}$ 273 323 373 $A,~\mathrm{\mu m}\cdot\text{K}$ 250.0 276.7 325.7 $\varkappa_\parallel,~\cfrac{\text{W}}{\text{m}\cdot\text{K}}$ 37.7 34.1 29.0 \[\rho=1.000\cdot10^3\\n=1.02\]
\end{solcard}

\begin{solcard}{B5}{Using the least squares method, obtain a formula for finding the coefficient $\Delta T_{\max}$ for each temperature value $T_0$ and for each material. Express your answer in terms of sample average values of products of the form $\Delta T^{k_1}d^{k_2}$, where $k_{1,2}$ are some numbers.}{0.60}

According to the method described in the condition, we compose a system of equations:\[\begin{cases} \cfrac{\partial S}{\partial\left( \Delta T_{\max}\right)}=0\\ \cfrac{\partial S}{\partial C}=0\end{cases}, \quad  \text{where} \quad\quad S=\sum_i\left(\Delta T_{i\ \mathrm{exp}}-\Delta T_{i\ \mathrm{th}}(d_i, A, B)\right)^2=\sum_i\left(\Delta T_i-\Delta T_{\max}+Cd_i\right)^2\]Thus, the system will take the form:\[\begin{cases} \cfrac\partial{\partial \left(\Delta T_{\max}\right)}\sum_i\left(\Delta T_i-\Delta T_{\max}+Cd_i\right)^2=0\\\cfrac\partial{\partial C}\sum_i\left(\Delta T_i-\Delta T_{\max}+Cd_i\right)^2=0\end{cases}\implies\begin{cases}0=\sum_i\left(\Delta T_i-\Delta T_{\max}+Cd_i\right)\\0=\sum_i\left(\Delta T_id_i-\Delta T_{\max} d_i+Cd_i^2\right)\end{cases}\]This is a system of linear equations for the coefficients $A$ and $B$. The coefficients in it are necessary products of the form $\Delta T^{k_1}d^{k_2}$. Using the average values of these products over the experimental sample, this system is written somewhat more elegantly in the form:\[\begin{cases}0=\overline{\Delta T}-\Delta T_{\max}+C\overline d\\0=\overline{\Delta Td}-\Delta T_{\max}\overline d+C\overline{d^2}\end{cases}\]Solving this system for $\Delta T_{\max}$, we finally obtain:\[\Delta T_{\max}=\frac{\overline{\Delta T}\cdot\overline{d^2}-\overline d\cdot\overline{\Delta Td}}{\overline{d^2}-\overline d^2}.\]

\par\smallskip\noindent\textbf{Answer:} \[\Delta T_{\max}=\frac{\overline{\Delta T}\cdot\overline{d^2}-\overline d\cdot\overline{\Delta Td}}{\overline{d^2}-\overline d^2}\]
\end{solcard}

\begin{solcard}{B6}{For each of the materials, find $\Delta T_{\max}$ for each of the temperatures $T_0$.}{1.20}

You can, of course, use the formula from the previous paragraph, but since the modeling function is simply a linear relationship, the value $\Delta T_{\max}$ can be calculated using least squares and using a standard calculator.

\par\smallskip\noindent\textbf{Answer:} $T_0,~\mathrm{K}$ 273 323 373 $\overset{\mathrm{MoS_2:}}{\Delta T_{\max},~\mathrm{K}}$ 2.841 2.451 2.159 $\overset{\mathrm{WS_2:}}{\Delta T_{\max},~\mathrm{K}}$ 3.155 2.679 2.293
\end{solcard}

\begin{solcard}{B7}{For each of the materials, find $m$. What makes the main contribution to heat transfer between the sample and the environment – thermal conductivity or radiation? Underline the correct option on your answer sheet.}{0.80}

You can calculate coefficients $\Delta T_{\max,0}$ and $m$ using least squares for a modeling function of the form: \[\Delta T_{\max}=\frac{\Delta T_{\max,0}}{T_0^m}\] using a standard calculator. As a result, we obtain: for $\rm MoS_2$, $m=0.88$ – the main contribution is made by $\underline{thermal conductivity}$ / radiation; for $\rm WS_2$, $m=0.83$ – the main contribution is made by $\underline{thermal conductivity}$ / radiation.


You can calculate coefficients $\Delta T_{\max,0}$ and $m$ using least squares for a modeling function of the form: \[\Delta T_{\max}=\frac{\Delta T_{\max,0}}{T_0^m}\] using a standard calculator. As a result we get:


The main contribution to heat transfer between the sample and the environment is thermal conductivity.

\par\smallskip\noindent\textbf{Answer:} Material $\rm MoS_2\rm WS_2m$ 0.88 0.83 Thermal conductivity makes the main contribution to the heat transfer between the sample and the environment.
\end{solcard}

\end{document}